\section{实直线上的紧致子空间}\label{section:ch03-compact-subspaces-real}

上一节中的某些定理使我们能够从已知的紧致空间来构造新的紧致空间，但是要想得到更多的紧致空间，我们必须先有一些紧致空间。最自然的办法当然还是从实直线开始，我们将证明 $\mathbb{R}$ 中每一个闭区间都是紧致的。其应用包括微积分中的极值定理和一致连续性定理的适当形式的推广。我们也将对 $\mathbb{R}^{n}$ 中的所有紧致空间给出一个刻画，也将给出实数集不可数的证明。

为了证明 $\mathbb{R}$ 中每一个闭区间是紧致的，我们仅需要用到实直线的序性质中的一个条件，即上确界性质。只用这一个性质便可以证明这个定理。因此，这个定理不仅适用于实直线，也同样适用于良序集和其他全序集。

\begin{theorem}\label{theorem:ch03-lots-interval-compact}
设 $X$ 是具有上确界性质的一个全序集。则关于序拓扑，$X$ 中的每一个闭区间都是紧致的。
\end{theorem}

\begin{proof}
第一步。给定 $a<b$，设 $\mathcal{A}$ 是 $[a,b]$ 的一个覆盖，它的成员是 $[a,b]$ 中关于子空间拓扑（与序拓扑相同）的开集。下面证明存在一个 $\mathcal{A}$ 的有限子族覆盖 $[a,b]$。首先证明：若 $x$ 是 $[a,b]$ 中异于 $b$ 的点，则在 $[a,b]$ 中存在点 $y>x$，使得区间 $[x,y]$ 可由 $\mathcal{A}$ 中最多两个成员覆盖。

若 $x$ 在 $X$ 中有直接后元，则令 $y$ 为这个直接后元。那么 $[x,y]$ 由两个点 $x$ 和 $y$ 所组成，所以它可为 $\mathcal{A}$ 中最多两个成员所覆盖。若 $x$ 在 $X$ 中没有直接后元，则选取 $\mathcal{A}$ 中包含 $x$ 的一个成员 $A$。由于 $x\neq b$ 以及 $A$ 是一个开集，所以对于 $[a,b]$ 中的某一个点 $c$ 有 $A$ 包含着形如 $[x,c)$ 的区间。在 $(x,c)$ 中选取一点 $y$，则区间 $[x,y]$ 被 $\mathcal{A}$ 中一个成员 $A$ 所覆盖。

第二步。设 $C$ 为 $[a,b]$ 中所有具有以下性质的点 $y>a$ 的集合：区间 $[a,y]$ 能够为 $\mathcal{A}$ 中有限多个成员所覆盖。对于 $x=a$，应用第一步的结论可见至少有一个这样的 $y$ 存在，从而 $C$ 是非空的。令 $c$ 是集合 $C$ 的上确界，则 $a<c\leqslant b$。

第三步。下面证明 $c$ 属于 $C$。也就是证明区间 $[a,c]$ 被 $\mathcal{A}$ 中有限个成员所覆盖。选取 $\mathcal{A}$ 中包含 $c$ 的一个成员 $A$。因为 $A$ 是开的，所以对于 $[a,b]$ 中某一个 $d$，$A$ 包含区间 $(d,c]$。若 $c$ 不属于 $C$，那么必然存在 $C$ 中一点 $z$ 包含在区间 $(d,c)$ 中（否则，$d$ 将是 $C$ 上比 $c$ 还要小的上界），参见\cref{figure:ch03-lots-compact-1}。由于 $z$ 在 $C$ 中，所以区间 $[a,z]$ 能被 $\mathcal{A}$ 中有限多个（比如说 $n$ 个）成员所覆盖。由于 $[z,c]$ 是包含在 $\mathcal{A}$ 的一个成员 $A$ 中，因此 $[a,c]=[a,z]\cup[z,c]$ 被 $\mathcal{A}$ 中 $n+1$ 个成员所覆盖。于是 $c$ 在 $C$ 中。这与假设矛盾。

第四步。最后，证明 $c=b$，从而完成定理的证明。假设 $c<b$，对于 $x=c$ 应用第一步得到：存在 $[a,b]$ 的一个点 $y>c$，使得区间 $[c,y]$ 能被 $\mathcal{A}$ 中有限个成员所覆盖，参见\cref{figure:ch03-lots-compact-2}。我们在第三步中已经证明了 $c$ 属于 $C$，从而 $[a,c]$ 可以被 $\mathcal{A}$ 中有限个成员所覆盖。因此区间
\begin{align*}
[a,y]=[a,c]\cup[c,y]
\end{align*}
也能被 $\mathcal{A}$ 中有限个成员所覆盖，这就意味着 $y$ 属于 $C$，与 $c$ 是 $C$ 的一个上界矛盾。
\end{proof}

\par
\begin{figure}[htbp]
\centering
\includegraphics[width=0.4\textwidth]{ch03-fig27-1.png}
\caption{第三步}\label{figure:ch03-lots-compact-1}
\end{figure}

\par
\begin{figure}[htbp]
\centering
\includegraphics[width=0.4\textwidth]{ch03-fig27-2.png}
\caption{第四步}\label{figure:ch03-lots-compact-2}
\end{figure}

\begin{corollary}\label{corollary:ch03-real-interval-compact}
$\mathbb{R}$ 中任何一个闭区间都是紧致的。
\end{corollary}

现在我们来刻划 $\mathbb{R}^{n}$ 中的紧致子空间：

\begin{theorem}\label{theorem:ch03-rn-compact-bounded-closed}
$\mathbb{R}^{n}$ 中一个子集 $A$ 是紧致的，当且仅当它是闭的并且就欧氏度量 $d$ 或平方度量 $\rho$ 而言是有界的。
\end{theorem}

\begin{proof}
只要考虑度量 $\rho$ 就可以了，因为不等式
\begin{align*}
\rho(x,y)\leqslant d(x,y)\leqslant\sqrt{n}\,\rho(x,y)
\end{align*}
保证了 $A$ 在 $\rho$ 中有界当且仅当它在 $d$ 中有界。

假设 $A$ 是紧致的。根据\cref{theorem:ch03-compact-subspace-hausdorff-closed}可见它是闭的。考虑开集族
\begin{align*}
\{\,B_{\rho}(\boldsymbol{0},m)\mid m\in\mathbb{Z}_{+}\,\},
\end{align*}
其并为 $\mathbb{R}^{n}$。于是有一个有限子族覆盖 $A$。从而，对于某一个 $M$ 有 $A\subset B_{\rho}(\boldsymbol{0},M)$。于是，对于 $A$ 中任意两点 $x$ 和 $y$，我们有 $\rho(x,y)\leqslant 2M$，因此，$A$ 对于 $\rho$ 而言是有界的。

反之，假设 $A$ 是闭的并且关于 $\rho$ 是有界的。假设对于 $A$ 中任意一对点 $x$ 和 $y$ 有 $\rho(x,y)\leqslant N$。选取 $A$ 中点 $x_{0}$，并且令 $\rho(x_{0},\boldsymbol{0})=b$，于是对于 $A$ 中任意 $x$，由三角不等式可以推出 $\rho(x,\boldsymbol{0})\leqslant N+b$。如果令 $P=N+b$，则 $A$ 是紧致方体 $[-P,P]^{n}$ 的一个子集合，又由于 $A$ 是闭集，所以 $A$ 也是紧致的。
\end{proof}

学生们常将这个定理记忆为：度量空间中的紧致子空间族就是有界闭集族。这显然荒谬，因为什么样的集合是有界的这与它的度量有关，而什么样的集合是紧致的则只依赖于空间的拓扑。

\begin{example}\label{example:ch03-compact-bounded-closed-examples}
$\mathbb{R}^{n}$ 中的单位球面 $S^{n-1}$ 和闭的单位球体 $B^{n}$ 都是紧致的，因为它们都是有界的闭集。集合
\begin{align*}
A=\{\,x\times(1/x)\mid 0<x\leqslant 1\,\}
\end{align*}
在 $\mathbb{R}^{2}$ 中是闭的，但它不是紧致的，因为它不是有界的。集合
\begin{align*}
S=\{\,x\times(\sin(1/x))\mid 0<x\leqslant 1\,\}
\end{align*}
在 $\mathbb{R}^{2}$ 中是有界的，但它不是紧致的，因为它不是闭的。
\end{example}

现在我们证明微积分中极值定理的适当推广形式。

\begin{theorem}[极值定理]\label{theorem:ch03-extreme-value}
设 $f\colon X\to Y$ 是连续的，其中 $Y$ 是具有序拓扑的全序集。若 $X$ 是紧致的，则在 $X$ 中存在点 $c$ 和 $d$，使得对于所有的 $x\in X$ 有 $f(c)\leqslant f(x)\leqslant f(d)$。
\end{theorem}

微积分中的极值定理只是这个定理的一种特殊情形，这只要在定理中取 $X$ 为 $\mathbb{R}$ 的闭区间，取 $Y$ 为 $\mathbb{R}$ 即可。

\begin{proof}
由于 $f$ 连续并且 $X$ 是紧致的，所以集合 $A=f(X)$ 是紧致的。我们证明 $A$ 有一个最大元 $M$ 和一个最小元 $m$。因而，由于 $m$ 和 $M$ 属于 $A$，必存在 $X$ 中点 $c$ 和 $d$ 使 $m=f(c)$ 和 $M=f(d)$。

若 $A$ 没有最大元，那么集合族
\begin{align*}
\{\,(-\infty,a)\mid a\in A\,\}
\end{align*}
就是 $A$ 的一个开覆盖。由于 $A$ 是紧致的，就有有限的子族
\begin{align*}
\{\,(-\infty,a_{1}),\cdots,(-\infty,a_{n})\,\}
\end{align*}
覆盖 $A$。设 $a_{i}$ 是 $a_{1},\cdots,a_{n}$ 中的最大者，则 $a_{i}$ 不属于这些集合中的任何一个，这与它们覆盖 $A$ 矛盾。

可以相仿地证明 $A$ 有最小元。
\end{proof}

现在我们来证明微积分中的一致连续性定理。为此，我们需要引入一个极为有用的概念，即度量空间开覆盖的 Lebesgue 数。作为准备，先给出以下概念：

\begin{definition}\label{definition:ch03-distance}
设 $(X,d)$ 是一个度量空间。$A$ 是 $X$ 的一个非空子集。对于每一个 $x\in X$，$x$ 到 $A$ 的\textbf{距离}（distance from $x$ to $A$）定义为
\begin{align*}
d(x,A)=\inf\{\,d(x,a)\mid a\in A\,\}.
\end{align*}
\end{definition}

易见，对于固定的 $A$，函数 $d(x,A)$ 是关于 $x$ 的连续函数：对于给定的 $x$ 和 $y$ 以及每一个 $a\in A$，不等式
\begin{align*}
d(x,A)\leqslant d(x,a)\leqslant d(x,y)+d(y,a)
\end{align*}
成立。从而，
\begin{align*}
d(x,A)-d(x,y)\leqslant\inf d(y,a)=d(y,A),
\end{align*}
因此，
\begin{align*}
d(x,A)-d(y,A)\leqslant d(x,y).
\end{align*}
交换 $x$ 和 $y$ 的位置，以上不等式仍然成立。由此推出 $d(x,A)$ 的连续性。

现在引进 Lebesgue 数的概念。前面讲过，度量空间中的一个有界子集 $A$ 的直径是指
\begin{align*}
\sup\{\,d(a_{1},a_{2})\mid a_{1},a_{2}\in A\,\}.
\end{align*}

\begin{lemma}[Lebesgue 数引理]\label{lemma:ch03-lebesgue-number}
设 $\mathcal{A}$ 为度量空间 $(X,d)$ 的一个开覆盖。若 $X$ 是紧致的，则存在 $\delta>0$ 使得 $X$ 的每一个直径小于 $\delta$ 的子集包含在 $\mathcal{A}$ 的某一元素之中。

数 $\delta$ 称为开覆盖 $\mathcal{A}$ 的一个\textbf{Lebesgue 数}（Lebesgue number）。
\end{lemma}

\begin{proof}
设 $\mathcal{A}$ 是 $X$ 的一个开覆盖。若 $X$ 本身是 $\mathcal{A}$ 的一个元素，那么任何一个正数都是 $\mathcal{A}$ 的 Lebesgue 数。以下假定 $X$ 不是 $\mathcal{A}$ 的元素。

选取 $\mathcal{A}$ 的子集的一个有限族 $\{A_{1},\cdots,A_{n}\}$ 覆盖 $X$。对于每一个 $i$，记 $C_{i}=X\setminus A_{i}$，通过取 $f(x)$ 为 $d(x,C_{i})$ 的平均数，定义一个映射 $f\colon X\to\mathbb{R}$。即
\begin{align*}
f(x)=\frac{1}{n}\sum_{i=1}^{n}d(x,C_{i}).
\end{align*}
我们证明：对每一个 $x$，$f(x)>0$。任意给定 $x\in X$，取 $i$ 使得 $x\in A_{i}$。再选取 $\varepsilon$ 使得 $x$ 的 $\varepsilon$-邻域包含于 $A_{i}$。则 $d(x,C_{i})\geqslant\varepsilon$，从而 $f(x)\geqslant\varepsilon/n$。

因为 $f$ 是连续的，它有极小值 $\delta$。我们证明 $\delta$ 就是所求的 Lebesgue 数。设 $B$ 为 $X$ 中一个直径小于 $\delta$ 的子集。在 $B$ 中取一点 $x_{0}$，则 $B$ 位于 $x_{0}$ 的 $\delta$-邻域中。那么
\begin{align*}
\delta\leqslant f(x_{0})\leqslant d(x_{0},C_{m}),
\end{align*}
其中 $d(x_{0},C_{m})$ 是所有 $d(x_{0},C_{i})$ 中的最大者。那么 $x_{0}$ 的 $\delta$-邻域被包含在开覆盖 $\mathcal{A}$ 的元素 $A_{m}=X\setminus C_{m}$ 之中。
\end{proof}

\begin{definition}\label{definition:ch03-uniformly-continuous}
设 $f$ 是从度量空间 $(X,d_{X})$ 到度量空间 $(Y,d_{Y})$ 的一个函数。若对任意 $\varepsilon>0$，存在 $\delta>0$ 使得对于 $X$ 的任何两点 $x_{0},x_{1}$，有
\begin{align*}
d_{X}(x_{0},x_{1})<\delta\Longrightarrow d_{Y}(f(x_{0}),f(x_{1}))<\varepsilon,
\end{align*}
则称函数 $f$ 是\textbf{一致连续的}（uniformly continuous）。
\end{definition}

\begin{theorem}[一致连续性定理]\label{theorem:ch03-uniform-continuity}
设 $f\colon X\to Y$ 是从紧致度量空间 $(X,d_{X})$ 到度量空间 $(Y,d_{Y})$ 的连续映射。则 $f$ 是一致连续的。
\end{theorem}

\begin{proof}
任意给定 $\varepsilon>0$，考虑由半径为 $\varepsilon/2$ 的球 $B(y,\varepsilon/2)$ 所组成的 $Y$ 的一个开覆盖。令 $\mathcal{A}$ 为这些球在 $f$ 下的原像所构成的 $X$ 的那个开覆盖。记 $\delta$ 为开覆盖 $\mathcal{A}$ 的一个 Lebesgue 数。则对于 $X$ 中满足 $d_{X}(x_{1},x_{2})<\delta$ 的两点 $x_{1},x_{2}$，由于这两点组成的集合 $\{x_{1},x_{2}\}$ 的直径小于 $\delta$，其像集 $\{f(x_{1}),f(x_{2})\}$ 必含在某一个开球 $B(y,\varepsilon/2)$ 之中。从而 $d_{Y}(f(x_{1}),f(x_{2}))<\varepsilon$。
\end{proof}

最后，我们来证明实数是不可数的。有趣的是，这个证明完全不用代数的方法（既不用十进位小数，也不用二进位小数或者其他什么方法），而仅仅用到 $\mathbb{R}$ 的序性质。

\begin{definition}\label{definition:ch03-isolated-point}
设 $X$ 是一个空间，$x\in X$。若单点集 $\{x\}$ 在 $X$ 中是开的，则称 $x$ 为 $X$ 的一个\textbf{孤立点}（isolated point）。
\end{definition}

\begin{theorem}\label{theorem:ch03-compact-hausdorff-uncountable}
设 $X$ 是一个非空的紧致的 Hausdorff 空间。若 $X$ 中没有孤立点，则 $X$ 是不可数的。
\end{theorem}

\begin{proof}
第一步。首先证明：给定 $X$ 的一个非空的开集 $U$ 和一点 $x\in X$，则存在一个包含于 $U$ 的非空开集 $V$，使得 $x\notin\overline{V}$。

我们总可以在 $U$ 中选取一个异于 $x$ 的点 $y$。当 $x$ 在 $U$ 中时，由于 $x$ 不是孤立点，可见 $y$ 的存在性。当 $x$ 不在 $U$ 中时，由于 $U$ 非空，也可见 $y$ 的存在性。然后分别选取包含 $x$ 和 $y$ 的无交开集 $W_{1}$ 和 $W_{2}$。那么集合 $V=U\cap W_{2}$ 便是我们所需的开集。$V$ 包含在 $U$ 中，它含有点 $y$，所以不是空集，并且 $V$ 的闭包不包含 $x$。参见\cref{figure:ch03-uncountable}。

第二步。我们证明：任何一个函数 $f\colon\mathbb{Z}_{+}\to X$，$f$ 都不是满射。由此可见，$X$ 是不可数的。

设 $x_{n}=f(n)$。对于非空的开集 $U=X$，应用第一步的结论，选取非空开集 $V_{1}\subset X$ 使得 $\overline{V}_{1}$ 不包含 $x_{1}$。一般地，对给定的非空开集 $V_{n-1}$，能够选取非空开集 $V_{n}$，使得 $V_{n}\subset V_{n-1}$ 并且 $\overline{V}_{n}$ 不包含 $x_{n}$。考虑 $X$ 的非空闭集套序列：
\begin{align*}
\overline{V}_{1}\supset\overline{V}_{2}\supset\cdots.
\end{align*}
因为 $X$ 是紧致的，根据\cref{theorem:ch03-fip-compact}，存在点 $x\in\bigcap\overline{V}_{n}$。因为 $x$ 属于 $\overline{V}_{n}$ 而 $x_{n}$ 不属于 $\overline{V}_{n}$，所以对于任意 $n$，点 $x$ 不等于 $x_{n}$。
\end{proof}

\par
\begin{figure}[htbp]
\centering
\includegraphics[width=0.4\textwidth]{ch03-fig27-3.png}
\caption{证明紧致 Hausdorff 空间不可数}\label{figure:ch03-uncountable}
\end{figure}

\begin{corollary}\label{corollary:ch03-real-interval-uncountable}
$\mathbb{R}$ 中的每一个闭区间都是不可数的。
\end{corollary}

\subsection*{习题}

\begin{enumerate}
\item 证明：若 $X$ 是一个全序集，并且它的每一个闭区间都是紧致的，则 $X$ 具有上确界性质。
\item 设 $X$ 是具有度量 $d$ 的度量空间，$A\subset X$ 是一个非空子集。
\begin{enumerate}
\item 证明 $d(x,A)=0$ 当且仅当 $x\in\overline{A}$.
\item 证明：若 $A$ 是一个紧致空间，则对于某 $a\in A$，有 $d(x,A)=d(x,a)$.
\item $A$ 在 $X$ 中的 $\varepsilon$-邻域定义为
\begin{align*}
U(A,\varepsilon)=\{\,x\mid d(x,A)<\varepsilon\,\}.
\end{align*}
证明 $U(A,\varepsilon)$ 等于所有开球 $B_{d}(a,\varepsilon)$ 之并，其中 $a\in A$.
\item 设 $A$ 是一个紧致子空间，$U$ 是包含 $A$ 的一个开集。证明：$A$ 的某 $\varepsilon$-邻域含在 $U$ 中.
\item 当 $A$ 是一个闭集但不紧致时，(d) 小题中的结论不成立.
\end{enumerate}
\item $\mathbb{R}_{K}$ 表示 $\mathbb{R}$ 赋予 $K$-拓扑。
\begin{enumerate}
\item 证明 $[0,1]$ 作为 $\mathbb{R}_{K}$ 的子空间不是紧致的.
\item 证明 $\mathbb{R}_{K}$ 是连通的。[提示：$(-\infty,0)$ 和 $(0,+\infty)$ 作为 $\mathbb{R}_{K}$ 的子空间也继承了其通常拓扑。]
\item 证明 $\mathbb{R}_{K}$ 不是道路连通的.
\end{enumerate}
\item 证明：多于一点的连通度量空间是不可数的。
\item 设 $X$ 是一个紧致的 Hausdorff 空间，$\{A_{n}\}$ 是 $X$ 的闭集的一个可数族。证明：若每一个 $A_{n}$ 在 $X$ 中有空内部，则它们的并 $\bigcup A_{n}$ 在 $X$ 中有空内部。[提示：仿照\cref{theorem:ch03-compact-hausdorff-uncountable}的证明] 这是我们将在第 8 章中讨论的 Baire 范畴定理的一个特殊情形。
\item 设 $A_{0}$ 是 $\mathbb{R}$ 中的闭区间 $[0,1]$，$A_{1}$ 是从 $A_{0}$ 中去掉它的``中间三分之一''$\big(\frac{1}{3},\frac{2}{3}\big)$ 而得到的集合，$A_{2}$ 是从 $A_{1}$ 中去掉它的两个``中间三分之一''$\big(\frac{1}{9},\frac{2}{9}\big)$ 和 $\big(\frac{7}{9},\frac{8}{9}\big)$ 而得到的集合。一般来说，$A_{n}$ 定义为
\begin{align*}
A_{n}=A_{n-1}-\bigcup_{k=0}^{\infty}\Big(\frac{1+3k}{3^{n}},\frac{2+3k}{3^{n}}\Big).
\end{align*}
它们的交
\begin{align*}
C=\bigcap_{n\in\mathbb{Z}_{+}}A_{n}
\end{align*}
称为\textbf{Cantor 集}（Cantor set），它是 $[0,1]$ 的子空间。
\begin{enumerate}
\item 证明 $C$ 是完全不连通的.
\item 证明 $C$ 是紧致的.
\item 证明每一个集合 $A_{n}$ 是有限多个长度为 $1/3^{n}$ 的无交闭区间的并，然后证明这些区间的端点属于 $C$.
\item 证明 $C$ 中没有孤立点.
\item 结论：$C$ 是不可数的.
\end{enumerate}
\end{enumerate}